\section{RL Teacher Regularization and Domain Weights}
\label{app:rl_teacher_connection}

Our experiments distill RL-trained Qwen3-1.7B teachers into a
Qwen3-1.7B-Base student. Teacher training introduces another source of
domain weighting: its reward scale and KL penalty determine the signal
passed to the student. We derive this connection by treating each teacher
as the optimum of its KL-regularized RL objective and distilling its
complete response distribution with reverse KL.

\paragraph{The teacher's RL objective.}
During RL, teacher $d$ is rewarded for solving its task and penalized for
deviating from a fixed reference policy $\mu_d$. The student starts
distillation from $\pi_{\theta_0}$, which in our experiments is
Qwen3-1.7B-Base. Thus $\mu_d$ specifies the reference in the teacher's
training objective, while $\pi_{\theta_0}$ specifies the student's
starting checkpoint.
For each domain $d$, fix a prompt distribution $\mathcal D_d$, a
sequence reward $R_d(x,y)$, and a KL penalty coefficient $\beta_d>0$.
Define
\begin{equation}
\begin{aligned}
\mathcal R_d(\theta)
&=\mathbb E_{x\sim\mathcal D_d,\,y\sim\pi_\theta(\cdot\mid x)}
  [R_d(x,y)],\\
K_d(\theta)
&=\mathbb E_{x\sim\mathcal D_d}
  D_{\mathrm{KL}}\!\left(\pi_\theta(\cdot\mid x)
  \|\mu_d(\cdot\mid x)\right),
\qquad J_d(\theta)=\mathcal R_d(\theta)-\beta_d K_d(\theta).
\end{aligned}
\label{eq:rl_teacher_objective}
\end{equation}
With common support and finite expectations and normalizers, the optimum
over response distributions has the standard form
\citep[Appendix A.1]{rafailov2023dpo}:
\begin{equation}
q_d^*(y\mid x)
=\frac{\mu_d(y\mid x)\exp(R_d(x,y)/\beta_d)}{Z_d(x)},
\qquad
Z_d(x)=\sum_y\mu_d(y\mid x)\exp(R_d(x,y)/\beta_d).
\label{eq:rl_teacher_optimum}
\end{equation}

\paragraph{Gradient equivalence.}
Consider the complete-sequence reverse KL, without response-length
normalization,
$L_d^{\mathrm{seq}}(\theta)=\mathbb E_{x\sim\mathcal D_d}
D_{\mathrm{KL}}(\pi_\theta(\cdot\mid x)\|q_d^*(\cdot\mid x))$.
Substituting Equation~\ref{eq:rl_teacher_optimum} gives
\begin{equation}
\begin{aligned}
L_d^{\mathrm{seq}}(\theta)
&=K_d(\theta)-\frac{\mathcal R_d(\theta)}{\beta_d}+C_d
=-\frac{J_d(\theta)}{\beta_d}+C_d,\\
-\nabla_\theta L_d^{\mathrm{seq}}(\theta)
&=\frac{1}{\beta_d}\nabla_\theta J_d(\theta),
\qquad C_d=\mathbb E_{x\sim\mathcal D_d}\log Z_d(x).
\end{aligned}
\label{eq:rl_teacher_gradient}
\end{equation}
Since $C_d$ is constant with respect to $\theta$, distilling this teacher
has the same gradient direction as maximizing its RL objective, scaled
by $1/\beta_d$. The derivative includes the change in the student's
response distribution as its parameters change. The equality holds
throughout student training, for any initialization.

\paragraph{How teacher training changes domain weights.}
With domain weights $\lambda_d$, the mixed descent direction separates
into a reward term and a reference term:
\begin{equation}
-\nabla_\theta\sum_d\lambda_d L_d^{\mathrm{seq}}(\theta)
=\sum_d\frac{\lambda_d}{\beta_d}\nabla_\theta\mathcal R_d(\theta)
  -\sum_d\lambda_d\nabla_\theta K_d(\theta).
\label{eq:rl_teacher_mixture}
\end{equation}
The reward terms push the student toward higher task rewards; the
reference terms pull it toward the teachers' RL references. The coefficient
of domain $d$'s reward gradient is therefore $\lambda_d/\beta_d$.
For example, with equal domain weights and comparable reward units,
mathematics and coding teachers with $\beta_d=0.05$ and $0.10$ give
the mathematics reward gradient twice the coefficient of the coding
reward gradient. The resulting update also depends on the magnitudes
and directions of these gradients and the reference terms.

If $R_d=s_d\widetilde R_d$ for $s_d>0$, the coefficient of
$\nabla_\theta\mathbb E[\widetilde R_d]$ becomes
$\lambda_d s_d/\beta_d$. Thus reward scale and KL strength act together:
rescaling both by the same positive constant leaves the optimal
teacher unchanged.

\paragraph{What student initialization changes.}
When the student starts from the same policy as every teacher's RL
reference, $\pi_{\theta_0}=\mu_d$, each KL term is minimized and
$\nabla_\theta K_d(\theta_0)=0$. The initial descent direction then
contains only the reward gradients:
\begin{equation}
\left.-\nabla_\theta\sum_d\lambda_d L_d^{\mathrm{seq}}(\theta)
\right|_{\theta_0}
=\left.\sum_d\frac{\lambda_d}{\beta_d}
  \nabla_\theta\mathcal R_d(\theta)\right|_{\theta_0}.
\label{eq:rl_teacher_initial_gradient}
\end{equation}
When the Base student's initial policy differs from a teacher's RL
reference, Equation~\ref{eq:rl_teacher_mixture} retains the corresponding
KL gradient from the first update. Distillation then combines improving
task reward with reducing the student's distance from that reference.

\paragraph{Connection to token balancing.}
The averaging rules in Section~\ref{sec:normalization} set the weights
assigned to domains and responses. The derivation above explains how
the teacher's RL objective sets the reward scale within a domain's
sequence loss. Our token losses in Section~\ref{sec:setting} use fixed
sampled prefixes, response-length averaging, and top-$k$ supervision
or clipped advantages. Their local gradients depend on these loss
choices as well as the teacher distribution; Equation~\ref{eq:rl_teacher_gradient}
describes the full derivative of the complete-sequence objective.
